\section{Related Work}\label{sec:related-work}

\subsection{Geospatial question answering and benchmarks}
Geospatial question answering requires interpreting geographic entities and relations in addition to matching textual content. MapQA provides question-answer pairs linked to geospatial entities and geometries, and includes question types involving neighborhood inference and entity types~\cite{li2025mapqa}. Its analysis highlights a distinction central to our setting: embedding-based retrieval can be useful for semantic candidate discovery, whereas questions requiring explicit spatial computation or multiple reasoning steps demand additional mechanisms. More recent benchmarks broaden the task and evaluation scope. GS-QA includes diverse spatial predicates, answer types, and questions that combine OpenStreetMap with other sources~\cite{saeedan2026gsqa}. These benchmarks provide complementary test beds; they do not by themselves prescribe how a system should represent intermediate dependencies or preserve evidence when a branch is unresolved.

\subsection{Spatial retrieval and GeoSPARQL}
Spatial-RAG combines sparse spatial-database retrieval with dense semantic retrieval and uses a multi-objective ranking strategy for spatial questions~\cite{yu2025spatialrag}. Agentic Tree-RAG shares the motivation to combine spatial and semantic evidence, while emphasizing explicit task decomposition, conditional execution, and provenance across intermediate steps. This distinction is architectural: Spatial-RAG describes a retrieval-and-generation framework, whereas our query tree is an executable investigation plan with declared producers, consumers, and outcome routes. A comparative performance claim requires evaluation under matched datasets and metrics.

GeoSPARQL standardizes RDF vocabularies for geographic features and geometries and extends SPARQL with spatial functions and relations~\cite{ogc2023geosparql}. It therefore offers a principled interface for spatially enabled knowledge graphs. In our design, GeoSPARQL is a spatial tool for graphs with verified geometry bindings; it is not a substitute for retrieving source records, establishing category membership, or resolving identity across datasets. This keeps ontology semantics distinct from observed instance-level evidence.

\subsection{Reasoning, acting, and structured plans}
ReAct interleaves language-model reasoning with actions that query external sources or interact with an environment~\cite{yao2023react}. Tree of Thoughts explores alternative reasoning paths and uses search over intermediate textual thoughts~\cite{yao2023tot}. These approaches establish useful patterns for action and branching, but their trees are not necessarily dataflow graphs whose nodes declare evidence inputs and whose transitions govern tool execution. Agentic Tree-RAG uses a query-specific tree as an operational contract: the planner declares tasks and routes, tools return evidence, and a dedicated manager validates state updates. This explicit separation is intended to make prerequisites, unresolved paths, and evidence provenance inspectable.

\subsection{Geospatial workflow extraction}
GeoAutoModuler studies structured extraction of functional modules, inputs, outputs, and dependencies from heterogeneous geospatial modeling code, including a staged skeleton-of-thought prompting strategy~\cite{liang2026geoautomoduler}. Its focus is understanding existing geospatial scripts and workflows. Our focus is different: planning and executing a natural-language geographic investigation over data sources at query time. The shared concern with explicit modules and dependencies is a useful design analogy, but GeoAutoModuler does not provide evidence for the effectiveness of our runtime planner or benchmark results.
